% from ega3-1-fr.tex
\oldpage[0\textsubscript{III}]{77}
Il résulte de (\hyperref[0.13.7.4-fr]{13.7.4}) appliqué au cas où $C$
est la catégorie des faisceaux de groupes abéliens sur $X$, $C'$ la
catégorie des groupes abéliens, et $\mathrm T=\Gamma$, que l'on a un
isomorphisme canonique
$\operatorname{gr}^p(\mathrm R^{n+1}\Gamma(\mathcal F))
\xrightarrow{\sim}\mathrm E_\infty^{p,n-p+1}(\mathcal F)$ pour tout
$p\in\mathbb Z$. D'autre part, comme en vertu de
(\hyperref[0.13.7.4-fr]{13.7.4}), le système projectif
$(H^n(X,\mathcal F_k))_{k\in\mathbb Z}$ vérifie (ML), on déduit de
(\hyperref[0.13.3.1-fr]{13.3.1}) un isomorphisme canonique
\begin{equation}
\label{0.13.7.5.1b-fr}
\tag{13.7.5.1}
H^{n+1}(X,\mathcal F)\xrightarrow{\sim}
\varprojlim_k H^{n+1}(X,\mathcal F_k).
\end{equation}

Comme le système projectif $\mathrm R^{n+1}\Gamma(\mathcal F)$ vérifie
(ML) en vertu de (\hyperref[0.13.7.4-fr]{13.7.4}), on a un isomorphisme
canonique
\[
\operatorname{gr}^p(\mathrm R^{n+1}\Gamma(\mathcal F))
\xrightarrow{\sim}
\varprojlim_k\operatorname{gr}^p(H^{n+1}(X,\mathcal F_k))
\]
(\hyperref[0.13.4.3-fr]{13.4.3}), et un isomorphisme canonique
\[
\varprojlim_k\operatorname{gr}^p(H^{n+1}(X,\mathcal F_k))
\xrightarrow{\sim}
\operatorname{gr}^p(\varprojlim_k H^{n+1}(X,\mathcal F_k))
\]
(\hyperref[0.13.4.5-fr]{13.4.5}). Tout revient donc à voir que
l'isomorphisme (\hyperref[0.13.7.5.1-fr]{13.7.5.1}) est compatible avec
les filtrations des deux membres~; mais cela résulte aussitôt des
définitions et de la commutativité du diagramme
\[
\xymatrix{
H^{n+1}(X,\mathcal F)\ar[r]^{\sim}\ar[dr]
&\varprojlim_k H^{n+1}(X,\mathcal F_k)\ar[d]\\
&H^{n+1}(X,\mathcal F_{p-1})
}
\]
pour tout $p$.
\end{proof}

\begin{env}[13.7.6]
\label{0.13.7.6-fr}
Soit $S$ un anneau muni d'une filtration
$(\mathrm F^i(S))_{i\in\mathbb Z}$ telle que $\mathrm F^0(S)=S$ (donc
$\operatorname{gr}^i(S)=0$ pour $i<0$). Supposons que chacun des $A_k$,
muni de la filtration définie dans
(\hyperref[0.13.7.1-fr]{13.7.1}), soit un $S$-$C$-module filtré
(\hyperref[0.13.6.6-fr]{13.6.6}), les morphismes $A_h\to A_k$ pour
$k\leq h$ étant des morphismes pour la structure de $S$-$C$-module
filtré~; nous dirons pour abréger que $\mathbf A$ est un \emph{système
projectif de $S$-$C$-modules filtrés}. Alors il est immédiat que les
morphismes $\mathrm B_r^{pq}(A_h)\to\mathrm B_r^{pq}(A_k)$ et
$\mathrm Z_r^{pq}(A_h)\to\mathrm Z_r^{pq}(A_k)$ pour $k\leq h$,
$1\leq r\leq+\infty$, sont des morphismes pour les structures de
$\operatorname{gr}^\bullet(S)$-$C'$-module bigradué
(\hyperref[0.13.6.5-fr]{13.6.5}), et que les familles
$(\mathrm Z_r^{pq}(\mathbf A))$, $(\mathrm B_r^{pq}(\mathbf A))$ et
$(\mathrm E_r^{pq}(\mathbf A))$ sont des
$\operatorname{gr}^\bullet(S)$-$C'$-modules bigradués pour $r$ fini,
les deux premiers étant des sous-modules de
$(\mathrm E_1^{pq}(\mathbf A))$. On notera encore
$\mathrm Z_r^{(n)}(\mathbf A)$, $\mathrm B_r^{(n)}(\mathbf A)$,
$\mathrm E_r^{(n)}(\mathbf A)$ les sous-objets respectifs des précédents
obtenus en ne prenant que les termes tels que $p+q=n$~; ce sont des
$\operatorname{gr}^\bullet(S)$-$C'$-modules gradués.

Lorsque le système $(\mathrm E_r^{pq}(\mathbf A))_{r\in\mathbb Z}$ est
essentiellement constant, $(\mathrm E_\infty^{pq}(\mathbf A))$ est donc
aussi un $\operatorname{gr}^\bullet(S)$-$C'$-module bigradué, et chaque
$\mathrm E_\infty^{(n)}(\mathbf A)$ un
$\operatorname{gr}^\bullet(S)$-$C'$-module gradué. En outre, les
$\beta^{p,n-p}:\mathrm E_\infty^{p,n-p}(A_k)\xrightarrow{\sim}
\operatorname{gr}^p(\mathrm R^n\mathrm T(A_k))$ constituent pour chaque
$k$ un isomorphisme pour la structure de
$\operatorname{gr}^\bullet(S)$-$C'$-module gradué de
$\mathrm E_\infty^{(n)}(A_k)$ sur
$\operatorname{gr}^\bullet(\mathrm R^n\mathrm T(A_k))$~; si on est dans
les conditions précédentes,
$\beta^{p,n-p}:\mathrm E_\infty^{(n)}(\mathbf A)\xrightarrow{\sim}
\varprojlim_k\operatorname{gr}^\bullet(\mathrm R^n\mathrm T(A_k))$
sera donc aussi un isomorphisme pour ces structures et il en est
évidemment de même de l'isomorphisme canonique
$\operatorname{gr}^\bullet(\mathrm R^n\mathrm T(\mathbf A))
\xrightarrow{\sim}
\varprojlim_k\operatorname{gr}^\bullet(\mathrm R^n\mathrm T(A_k))$,
donc les isomorphismes
(\hyperref[0.13.7.4.1-fr]{13.7.4.1}) constituent un isomorphisme pour les
structures de $\operatorname{gr}^\bullet(S)$-$C'$-module gradué.
\end{env}

\begin{proposition}[13.7.7]
\label{0.13.7.7-fr}
Soit $S$ un anneau noethérien $\mathfrak J$-adique. Supposons que $C$ soit
une catégorie abélienne dont tout objet est isomorphe à un sous-objet d'un
objet injectif, et soit $\mathrm T$ un foncteur covariant additif de $C$ dans
la catégorie des groupes abéliens. Soit
$\mathbf A=(A_k)_{k\in\mathbb Z}$ un
